% Frontespizio
\begin{titlepage}
\centering
{\Large Luigi E. Picasso\par}
\vspace{3cm}
{\Huge Lezioni\\[0.3em] di\\[0.3em] Fisica Generale 2\par}
\vfill
{\large EDIZIONI ETS\par}
\end{titlepage}

\renewcommand{\contentsname}{INDICE}
\tableofcontents
\clearpage

\section*{PREFAZIONE}

Il corso di Fisica Generale 2 tradizionalmente comprende l'elettromagnetismo classico e talvolta, come in questo caso, la teoria della relatività. Il problema è che quando si scrive un libro sull'elettromagnetismo si sa da dove cominciare ma non si sa dove finire: l'elettromagnetismo classico è uno dei capitoli più belli, ma anche uno dei più vasti della fisica classica.\\
Ma questo non è un trattato di elettromagnetismo, è solo un libro di lezioni tenute agli studenti del 2$^\circ$ anno e un corso ha una durata limitata e quindi impone necessariamente una scelta fra i tanti argomenti che noi docenti riteniamo che ``non possono essere ignorati''; in effetti il contenuto di questo libro va un po' oltre a quello che è normalmente il programma del corso.

Questo lavoro nasce dalla mia esperienza di insegnamento presso il Corso di Laurea in Matematica dell'Università di Pisa ed è probabile che tanto l'impostazione quanto la scelta degli argomenti siano stati condizionati dai ``gusti'' e dalla mentalità degli studenti di questa disciplina. L'impostazione non è quella ``storica'', secondo la quale l'elettromagnetismo ha preceduto, e anzi determinato, la teoria della relatività: sebbene venga anche brevemente esposto lo sviluppo storico dell'elettromagnetismo, sia l'introduzione del campo magnetico sia le equazioni di Maxwell si appoggiano in maniera determinante sul principio di relatività e sull'invarianza della carica elettrica. In questo senso si muove anche l'impostazione dell'ottimo libro sull'elettromagnetismo de ``La fisica di Berkeley'' e gli studenti sembrano apprezzare l'economia concettuale di questa impostazione.\\
Ogni docente ha le sue fissazioni e al lettore non ci vorrà molto tempo per scoprire che la mia è l'uso delle simmetrie: sono uno strumento molto potente e, specialmente nel caso del campo magnetico che ha proprietà di trasformazione non banali per riflessioni, permette di risolvere problemi che a prima vista si presentano come ``casi disperati''.

L'elenco delle omissioni sarebbe più lungo dell'indice di queste dispense, dato che ho cercato di privilegiare i concetti di base rispetto ai molti e pur importanti sviluppi dell'elettromagnetismo. Nel caso degli studenti di Fisica i corsi di laboratorio vengono in soccorso a integrare quegli argomenti che qui o non sono trattati (p.es.\ strumenti di misura, ottica geometrica) o dei quali vengono fornite solo le basi teoriche (circuiti, motori, ecc.).\\
In particolare il capitolo sui fenomeni ondulatori (interferenza, diffrazione) è molto ridotto: per ragioni di tempo, per l'impossibilità di utilizzare gli strumenti matematici necessari \dots; l'obiettivo non è quello di una discussione ragionevolmente esauriente di questi argomenti, ma quello di mettere in grado lo studente di seguire i ragionamenti che normalmente servono per introdurre i fondamenti della Meccanica Quantistica (interferenza con singolo fotone, esperimenti con il polaroid, ecc.): è auspicabile che al 3$^\circ$ o al 4$^\circ$ anno anche gli studenti di Matematica imparino in cosa consiste la seconda rivoluzione scientifica del XX$^\circ$ secolo.

Molti sono i colleghi che ringrazio per le loro osservazioni e consigli: le Dott.$^{\mathit{sse}}$ L.~Bosman e S.~Degl'Innocenti, i proff.\ L.~Bracci, R.~Mannella e in modo particolare il prof.\ P.~Menotti per le discussioni veramente illuminanti su un problema non proprio banale, quello che viene riportato nel capitolo 11.4. Sia chiaro che questi colleghi non sono responsabili degli eventuali errori nei quali posso essere incorso: nell'elettromagnetismo il rischio di dire cose inesatte è tutt'altro che trascurabile; anzi, direi che quando si insegna questa materia \dots\ c'è sempre da imparare.\\
Veramente con piacere ringrazio gli studenti che hanno seguito questo corso nell'anno accademico 1997-98, anno in cui ho iniziato la stesura di queste dispense, per il contributo positivo che molti di loro (studenti veramente molto bravi) con le loro domande ed il loro spirito critico, inconsapevolmente hanno dato a questo lavoro.

L'immagine riportata sul frontespizio è una riproduzione di una incisione di Paolo Lasinio ripresa dalla Tribuna di Galileo: ``Si presenta in essa quell'adunanza dell'istituto di Francia, nella quale il celebre Volta espose al Consesso degli uomini più dotti d'Europa le importanti esperienze della famosa sua Pila''. (Fra gli altri, sono riprodotti Napoleone, Monge, Laplace, Legendre e Biot).\\
Ringrazio il Prof.\ Roberto Vergara Caffarelli, responsabile del Centro per la conservazione e lo studio degli strumenti scientifici, e il personale del Centro, per avermi consigliato e fornito questa bella riproduzione.

\begin{flushright}
Luigi E. Picasso
\end{flushright}

Pisa, ottobre 1999
\clearpage
